\documentclass[10pt,a4paper]{amsart}
\usepackage[margin=19mm]{geometry}
\usepackage{amsmath,amssymb,mathtools}
\usepackage[scr=rsfs,cal=euler]{mathalfa}
\usepackage{xfrac}
\usepackage[unicode,hidelinks,bookmarks=false]{hyperref}
\newcommand{\ie}{i.e.\ }
\newcommand{\cf}{cf.\ }
\newcommand{\resp}{respectively\ }
\newcommand{\Subsection}{\subsection}
\providecommand{\nobreakdash}{\nobreak\hbox{-}}
\setlength{\emergencystretch}{2em}
\multlinegap0pt
\usepackage{amssymb}
\usepackage{tcolorbox}
\usepackage[shortlabels]{enumitem}
\usepackage{mathtools}
\newtheorem{lemma}{Lemma}
\newtheorem{theorem}{Theorem}
\newcommand{\R}{\mathbb{R}}
\title{A unified approach to the divergence equation\\ and related functional inequalities}
\author{Filippo Gazzola, Hans-Christoph Grunau, Gianmarco Sperone}
\date{Source: arXiv:2608.25861v1 (CC BY 4.0)}
\begin{document}
\maketitle

\noindent\emph{Source and licence.} A unified approach to the divergence equation\\ and related functional inequalities. Version arXiv:2608.25861v1 (CC BY 4.0), distributed by arXiv under the Creative Commons Attribution 4.0 International licence, http://creativecommons.org/licenses/by/4.0/, as recorded in the arXiv licence metadata for this exact version and shown on the abstract page https://arxiv.org/abs/2608.25861v1. Full licence notice: \texttt{licenses/arxiv-2608.25861-CC-BY-4.0.txt}.\par\smallskip

\noindent\textit{Editorial scope.} Counted unit: the article's theorem lowerbound (source0.tex lines 450--472). The article's complete original proof of it is delivered with it (source0.tex lines 473--533, one \texttt{proof} environment of 61 lines). No mathematics is written, repaired or re-derived: the statement and the proof are the article's own lines, in the article's own notation, and no retained line is substituted or re-flowed.  Retained uncounted as the source-local context the proof needs: lines 128--130; lines 331--340; lines 359--380; lines 366--368; lines 428--430.  Nothing was omitted from any retained span, so the package carries no omission record.  No retained line contains a citation, so the wrapper carries no bibliography and no bibliography entry.  No retained line contains a figure environment, an image-inclusion command or drawing code, so no image is delivered and the package declares no asset dependency.  Editorial formatting only, in addition: an AMS \texttt{amsart} preamble that restates the article's own declarations and macros used by the retained lines, namely the environments \texttt{lemma}, \texttt{theorem} on separate counters, together with the standard packages those lines need.  The retained environments print in this document as Lemma 1, Theorem 1 in order of appearance, whereas the article prints the same items with the numbering of its own full document.  The complete native source of the article as distributed in the arXiv e-print for this exact version is bundled unchanged as \texttt{sources/208600/source0.tex}.
\begin{equation}\label{Bogdom}
\nabla\cdot u=f \quad \mbox{in} \quad \Omega\, ,\qquad u=0 \quad \mbox{on} \quad \partial\Omega
\end{equation}

\begin{lemma}\label{lem:lemma1}
For any vector field
\begin{equation}\label{H10s}
v \in H^{1}_{0,\sigma}(\Omega,\R^n)\doteq\{w\in H^1_0(\Omega,\R^n);\, \nabla \cdot w=0\mbox{ in }\Omega\},
\end{equation}
any $j\in\{1,\ldots,n\}$ and any $g\in\mathcal{C}^1(\R,\R)$, one has
\begin{equation} \label{exlemma}
\int_\Omega g'(x_j )v_j(x)\, dx=0.
\end{equation}	
\end{lemma}

\begin{theorem}\label{minBog}
For any $f\in L^2_0(\Omega,\R)$ there exists a unique solution
$u_f\in H^1_0(\Omega,\R^n)$ to \eqref{Bogdom} such that
\begin{equation}\label{inf}
\|\nabla u_f\|_{L^2(\Omega)}^2=\inf\Big\{\|\nabla v\|_{L^2(\Omega)}^2\, ;\ v\in H^1_0(\Omega,\R^n)\, ,\ \nabla\cdot v=f\text{ in }\Omega\Big\}\, .
\end{equation}
Moreover, with $H^1_{0,\sigma}(\Omega,\R^n)$ as in \eqref{H10s}, satisfying the following Euler-Lagrange equation
\begin{equation}\label{ELeq}
\int_\Omega \nabla u : \nabla v\, dx=0\qquad\forall v\in H^1_{0,\sigma}(\Omega,\R^n)
\end{equation}
is necessary and sufficient for a solution $u\in H^1_0(\Omega,\R^n)$ of \eqref{Bogdom} to minimise \eqref{inf} so that $u\equiv u_f$. Equivalently,
there exists a unique $p_{f} \in L^2_0(\Omega,\R)$ such that the pair $(u_{f},p_{f})$ weakly solves the generalised Stokes system
\begin{equation} \label{gstokes}
\left\{
\begin{aligned}
	& - \Delta u_{f} + \nabla p_{f} = 0 \, , \quad \nabla \cdot u_{f} = f \ \ \mbox{ in } \ \ \Omega \, , \\[5pt]
	& u_{f}=0 \ \ \mbox{ on } \ \ \partial \Omega  \, .
\end{aligned}
\right.
\end{equation}
Moreover, if $\Omega$ has a boundary of class $\mathcal{C}^{m+1,1}$ and $f\in H^{m}(\Omega, \mathbb{R}) \cap L^2_0(\Omega, \mathbb{R})$, for some $m \in \mathbb{N}$, then $u_f\in H^{m+1}(\Omega, \mathbb{R}^{n}) \cap H^1_0(\Omega, \mathbb{R}^{n})$.
\end{theorem}

\begin{equation}
\int_\Omega \nabla u : \nabla v\, dx=0\qquad\forall v\in H^1_{0,\sigma}(\Omega,\R^n)
\end{equation}

\begin{equation}\label{bogo}
C_B(\Omega) \doteq \sup_{f\in L^2_0(\Omega,\R) \setminus \{0\}} \ \dfrac{\| \nabla u_{f} \|_{L^2(\Omega)}^2}{\| f \|_{L^2(\Omega)}^2} \, ,
\end{equation}

\begin{theorem}\label{lowerbound}
	For any bounded Lipschitz domain $\Omega\subset\R^n$ one has
	$$
	C_B(\Omega)\ge n \, .
	$$
Moreover, $C_B(\Omega)>n$ whenever one of the following two facts occurs:\\
$\bullet$ there exist $i \in \{1,...,n\}$ and $g\in\mathcal{K}_{i}(\Omega, \mathbb{R}) \cap L^{2}(\Omega, \mathbb{R})$ such that the unique solution $w \in H^{1}_{0}(\Omega, \mathbb{R})$ of
\begin{equation}\label{generalTorsion}
- \Delta w = g\mbox{ in }\Omega\, , \qquad w = 0\mbox{ on }\partial \Omega  \, ,
	\end{equation}
	satisfies
	\begin{equation}\label{condition}
	\int_\Omega \left| \dfrac{\partial w}{\partial x_{i}} \right|^2 \, dx<\frac{1}{n}\int_\Omega|\nabla w|^2\, dx\, ,
	\end{equation}
$\bullet$ there exists $h\in \mathcal{C}(\R, \mathbb{R})$ such that the unique solution to
\begin{equation}\label{hsum}
-\Delta w=h\left(\sum_{i=1}^nx_i\right)\mbox{ in }\Omega\, ,\qquad w=0\mbox{ on }\partial\Omega\, ,
\end{equation}
satisfies
\begin{equation}\label{negativepd}
\sum_{1=i<j}^{n}\int_\Omega\partial_{x_i}w\, \partial_{x_j}w\, dx<0\, .
\end{equation}
\end{theorem}

\noindent\begin{proof} We start by noticing the elementary inequality for scalar functions
	\begin{equation}\label{elementary}
	\int_\Omega|\nabla w|^2 \, dx=\sum_{i=1}^{n} \left( \int_\Omega \left| \dfrac{\partial w}{\partial x_{i}} \right|^2 \, dx \right) \ge n \min_{k \in \{1,...,n \} }\ \int_\Omega \left| \dfrac{\partial w}{\partial x_{k}} \right|^2\, dx \qquad
	\forall w\in H^1_0(\Omega, \mathbb{R})\, .
	\end{equation}
	We apply \eqref{elementary} to the weak solution $w\in H^1_0(\Omega, \mathbb{R})$ of the {\em torsion problem}
\begin{equation}\label{torsionpb}
- \Delta w = 1\mbox{ in }\Omega \, ,\qquad w = 0\mbox{ on }\partial \Omega \, ,
\end{equation}
that is,
\begin{equation}\label{partialk0}
\int_\Omega \nabla w \cdot \nabla \varphi \, dx = \int_\Omega \varphi \, dx\qquad \forall \varphi \in H^1_0(\Omega, \mathbb{R}) \, .
\end{equation}
Let  $k\in\{1,...,n\}$ denote  the index of the partial derivative of $w$ attaining the minimum in the right hand side of \eqref{elementary};
	clearly, there may be more than one such index, e.g.\ in domains with some symmetry property. Then, \eqref{elementary} becomes
	\begin{equation}\label{partialk}
	\int_\Omega|\nabla w|^2 \, dx\ge n\, \int_\Omega \left| \dfrac{\partial w}{\partial x_{k}} \right|^2 \, dx \, .
	\end{equation}
	Take the vector field $u=(u_1,...,u_n)\in H^1_0(\Omega, \mathbb{R}^{n})$ such that $u_k=w$ and $u_i\equiv0$ for $i\neq k$, and define $f \in L_{0}^{2}(\Omega, \R)$ by
	$$
	f\doteq\nabla\cdot u=\dfrac{\partial w}{\partial x_{k}} \quad \text{in} \ \ \Omega \, .
	$$
	We claim that $u = u_f$ as in Theorem \ref{minBog}; given any $v = (v_{1},...,v_{n}) \in H^{1}_{0,\sigma}(\Omega, \mathbb{R}^{n})$, notice that
	$$
	\int_{\Omega} \nabla u : \nabla v \, dx= \int_{\Omega} \nabla w \cdot \nabla v_{k}\, dx = \int_{\Omega} v_{k}\, dx = 0 \, ,
	$$
	owing to Lemma \ref{lem:lemma1} and \eqref{partialk0}, so that the Euler-Lagrange equation \eqref{ELeq} is fulfilled.
	From Theorem \ref{minBog} and \eqref{bogo}+\eqref{partialk} we then deduce
	$$
	C_B(\Omega)=\sup_ {h \in L^2_0(\Omega) \setminus \{0\}}\frac{\|\nabla u_h\|_{L^2(\Omega)}^2}{\|h\|_{L^2(\Omega)}^2} \geq \frac{\|\nabla u_f\|_{L^2(\Omega)}^2}{\|f\|_{L^2(\Omega)}^2}=\frac{\|\nabla w\|_{L^2(\Omega)}^2}{\|\partial_{x_k}w\|_{L^2(\Omega)}^2} \ge n \, ,
	$$
	thus providing the stated lower bound for $C_B(\Omega)$.\par
	In order to prove the first non-minimality criterion, consider $i \in \{1,...,n\}$, $g\in\mathcal{K}_{i}(\Omega, \mathbb{R}) \cap L^{2}(\Omega, \mathbb{R})$ and $w\in H^1_0(\Omega, \mathbb{R})$  satisfying \eqref{generalTorsion}-\eqref{condition}, meaning that
	\begin{equation}\label{partialk1}
	\int_\Omega \nabla w \cdot \nabla \varphi  \, dx = \int_\Omega g \, \varphi \, dx\qquad \forall \varphi \in H^1_0(\Omega, \mathbb{R}) \, .
	\end{equation}
	Take the vector field
	$u=(u_1,...,u_n)\in H^1_0(\Omega, \mathbb{R}^{n})$ such that $u_i=w$ and $u_j\equiv0$ for $j \neq i$, and define $h \in L_{0}^{2}(\Omega, \mathbb{R})$ by
	$$
	h\doteq\nabla\cdot u=\dfrac{\partial w}{\partial x_{i}} \quad \text{in} \ \ \Omega \, .
	$$
	As before, we claim that $u = u_h$; given any vector field $v = (v_{1},...,v_{n}) \in H^{1}_{0,\sigma}(\Omega, \mathbb{R}^{n})$, notice that
	$$
	\int_{\Omega} \nabla u : \nabla v \, dx= \int_{\Omega} \nabla w \cdot \nabla v_{i} \, dx= \int_{\Omega} g \, v_{i}\, dx = 0 \, ,
	$$
	owing to Lemma \ref{lem:lemma1} and \eqref{partialk1}, so that the Euler-Lagrange equation \eqref{ELeq} is fulfilled.  Then, by using \eqref{bogo} and \eqref{condition}, we find
	$$
	C_B(\Omega)\ge\frac{\|\nabla u_{h}\|_{L^2(\Omega)}^2}{\|h\|_{L^2(\Omega)}^2}=\frac{\|\nabla w\|_{L^2(\Omega)}^2}{\|\partial_{x_i}w\|_{L^2(\Omega)}^2}>n\, ,
	$$
	which completes the proof of the first criterion.\par
For the second criterion, let $w\in H^1_0(\Omega, \mathbb{R})$ satisfy both \eqref{hsum} and \eqref{negativepd}.
Take the vector field $u=(u_i,...,u_n)=(w,...,w)\in H^1_0(\Omega, \mathbb{R}^{n})$. Then \eqref{ELeq} is satisfied and
$$
|\nabla u|^2=n|\nabla w|^2\, ,\qquad|\nabla\cdot u|^2=|\nabla w|^2+2\sum_{1=i<j}^{n}\partial_{x_i}w\, \partial_{x_j}w\, .
$$
Therefore, in view of \eqref{negativepd}, we have
$$
\int_\Omega|\nabla u|^2\, dx=n\int_\Omega|\nabla\cdot u|^2\, dx-2n\sum_{1=i<j}^{n}\int_\Omega\partial_{x_i}w\, \partial_{x_j}w\, dx>n\int_\Omega|\nabla\cdot u|^2\, dx\, ,
$$
which proves again that $C_B(\Omega)>n$. We point out that this second criterion is equivalent to the first one after a careful
(and lengthy) change of variables. We provided here a direct proof.\end{proof}

\end{document}
